% Generated draft: requires mathematical and source review.
\documentclass[11pt]{article}
\usepackage[margin=25mm]{geometry}
\usepackage{fontspec}
\setmainfont{Latin Modern Roman}
\usepackage{amsmath,amssymb,mathtools,unicode-math}
\setmathfont{Latin Modern Math}
\usepackage{enumitem,needspace,xurl,hyperref}
\setlength{\parindent}{0pt}
\setlength{\parskip}{5pt}
\setlength{\emergencystretch}{3em}
\newcommand{\N}{\mathbb{N}}
\newcommand{\Z}{\mathbb{Z}}
\newcommand{\Q}{\mathbb{Q}}
\newcommand{\R}{\mathbb{R}}
\newcommand{\C}{\mathbb{C}}
\newcommand{\F}{\mathbb{F}}
\newcommand{\A}{\mathbb{A}}
\newcommand{\PP}{\mathbb{P}}
\newcommand{\E}{\mathbb{E}}
\newcommand{\eps}{\varepsilon}
\newcommand{\dd}{\,\mathrm{d}}
\begin{document}
\section*{UnsolvedMath ID: 30001289}
\textbf{Source problem code:} catalogue entry
\subsection{Definitions and mathematical statement}
The supplied question concerns the integrable Hamiltonian
\[
H_0(I)=\frac12\sum_{i=1}^{n}I_i^2.
\]
The number of action variables \(n\), their domain, the underlying symplectic phase space, and the differentiability index \(r\) are not specified in the record. Consequently, the following is a conditional formulation template, requiring these data to be supplied.

For a specified symplectic phase space \((M,\omega)\) on which the displayed \(H_0\) is defined, and a specified neighborhood \(\mathcal U\) of \(H_0\) in a specified \(C^r\) topology, ask whether there exists
\[
H\in C^\infty(M,\R)\cap\mathcal U
\]
whose unit-energy set
\[
\Sigma_H:=H^{-1}(\{1\})
\]
contains a dense Hamiltonian trajectory. More explicitly, define the Hamiltonian vector field by \(\iota_{X_H}\omega=-\dd H\). If \(\phi_H^t(x)\) denotes its maximal trajectory through \(x\), with maximal existence interval \(J_x\), the density requirement is
\[
\exists x\in\Sigma_H\qquad
\overline{\{\phi_H^t(x):t\in J_x\}}^{\,\Sigma_H}=\Sigma_H,
\]
where the closure is taken in the relative topology of the entire unit-energy set.

The source does not specify whether the requested trajectory must exist for all real times or whether density is required for a forward trajectory alone. It also does not state that \(1\) is a regular value of \(H\); thus calling \(\Sigma_H\) a surface does not supply a regularity hypothesis. Finally, if ``in a \(C^r\)-small neighborhood'' means arbitrarily close to \(H_0\), the intended quantifiers would be: for every neighborhood \(\mathcal U\) of \(H_0\), there exists such an \(H\). The record does not explicitly distinguish that interpretation from existence in a particular sufficiently small neighborhood.
\subsection{Short English statement}
The question asks whether a smooth Hamiltonian close to the quadratic integrable Hamiltonian can have one trajectory dense in its whole energy-one set. The phase space, dimension, regularity index, and precise meaning of closeness must be supplied before this becomes a fully specified mathematical problem. The supplied partial result concerns density in a positive-measure subset, which does not establish the requested density in the whole energy surface.
\subsection{Sources}
\begin{itemize}
\item Supplied local UnsolvedMath registry record, permanent ID \(30001289\), title: Dense Hamiltonian Trajectories Near Integrable Systems; source problem code: \texttt{catalogue entry}. The record gives proposed year \(2009\) and names no proposer.
\item Supplied source citation: \emph{Dynamische Systeme} (2010), Oberwolfach Reports, pp.\ 1787--1842, DOI \texttt{10.4171/owr/2009/32}. Source URL and literature-source link labeled ``OWR report'': \url{https://doi.org/10.4171/owr/2009/32}. The supplied background describes Oberwolfach Reports as documenting mathematical workshops and their open-problem discussions.
\item Supplied literature-source link labeled ``Herman question and partial theorem'': \url{https://ems.press/content/serial-article-files/46234}. No theorem number, author list, or theorem statement is supplied for this reference.
\item Literature assessment dated \texttt{2026-08-21}: the supplied curation labels the problem partially solved and reports that a cited theorem constructs an orbit dense in a positive-measure subset, but does not provide a trajectory dense in the whole energy surface. This is supplied dated curation derived from the OWR material, not an independently verified literature assessment; the source explicitly requests independent verification of time-sensitive status and cited claims.
\end{itemize}
\end{document}
